\documentclass[12pt]{article}
\usepackage{amsmath, amssymb, amsthm}
\usepackage{geometry}
\geometry{margin=1in}

\newtheorem{exercise}{Exercise}
\theoremstyle{remark}
\newtheorem*{solution}{Solution}
\newtheorem*{remark}{Remark}

\title{Random Walks --- Week 4 Piazza Exercises}
\author{CUNY, Spring 2026}
\date{}

\begin{document}
\maketitle

\section*{Week 4: Zero-One Laws and Limit Behavior}

% ----------------------------------------------------------------
\begin{exercise}[Limsup and liminf of events]
Let $(A_n)_{n \ge 1}$ be a sequence of events on a probability space $(\Omega, \mathcal{F}, P)$.
\begin{enumerate}
  \item[(a)] Show that $\liminf_{n\to\infty} A_n = \bigcup_{k=1}^\infty \bigcap_{n \ge k} A_n$
             and $\limsup_{n\to\infty} A_n = \bigcap_{k=1}^\infty \bigcup_{n \ge k} A_n$.
  \item[(b)] Show that $\liminf A_n \subseteq \limsup A_n$.
  \item[(c)] If $P(A_n \triangle A) \to 0$ as $n \to \infty$ for some event $A$,
             show that $P(A) = \lim_{n\to\infty} P(A_n)$ and
             $P(\liminf A_n) = P(\limsup A_n) = P(A)$.
\end{enumerate}
\end{exercise}

\begin{solution}
\textbf{(a)} By definition, $\omega \in \liminf A_n$ iff $\omega \in A_n$ for all but
finitely many $n$, i.e., there exists $k$ such that $\omega \in A_n$ for all $n \ge k$.
This is exactly $\omega \in \bigcup_k \bigcap_{n \ge k} A_n$.  Similarly,
$\omega \in \limsup A_n$ iff $\omega \in A_n$ for infinitely many $n$, which is
$\omega \in \bigcap_k \bigcup_{n \ge k} A_n$.

\textbf{(b)} If $\omega$ is in all $A_n$ eventually, it is in $A_n$ infinitely often.

\textbf{(c)} $|P(A_n) - P(A)| \le P(A_n \triangle A) \to 0$, so $P(A_n) \to P(A)$.
For the liminf: $\liminf A_n = \bigcup_k \bigcap_{n \ge k} A_n \uparrow$, and
$P(\bigcap_{n \ge k} A_n) \ge P(A) - P(A \triangle A_n) \to P(A)$, so
$P(\liminf A_n) = P(A)$. Then $P(\limsup A_n) \ge P(\liminf A_n) = P(A)$ and
$P(\limsup A_n) \le P(A) + P(A_n \triangle A) \to P(A)$, giving equality.
\end{solution}

% ----------------------------------------------------------------
\begin{exercise}[Tail $\sigma$-algebra and constants]
Let $X_1, X_2, \ldots$ be independent real-valued random variables.  Define the
\emph{tail $\sigma$-algebra}
\[
  \mathcal{T} = \bigcap_{n=1}^\infty \sigma(X_{n+1}, X_{n+2}, \ldots).
\]
Show that any $\mathcal{T}$-measurable random variable $Z$ must be $P$-a.s.\ equal to
a constant (i.e., there exists $c \in \mathbb{R}$ with $P(Z = c) = 1$).
\end{exercise}

\begin{solution}
By Kolmogorov's 0-1 Law, $\mathcal{T}$ is $P$-trivial: every $A \in \mathcal{T}$
satisfies $P(A) \in \{0,1\}$.  If $Z$ is $\mathcal{T}$-measurable, then for every
Borel set $B \subseteq \mathbb{R}$, the event $\{Z \in B\} \in \mathcal{T}$ has
probability $0$ or $1$.

Let $F(x) = P(Z \le x)$ be the distribution function of $Z$.  Since $F$ is
right-continuous, non-decreasing, with $F(-\infty)=0$ and $F(+\infty)=1$, and since
$\{Z \le x\} \in \mathcal{T}$ forces $F(x) \in \{0,1\}$ for every $x$, the function
$F$ must be the Heaviside step function at some point $c$:
\[
  c = \sup\{x : F(x) = 0\} = \inf\{x : F(x) = 1\}.
\]
Therefore $P(Z = c) = 1$. \hfill$\square$
\end{solution}

% ----------------------------------------------------------------
\begin{exercise}[Percolation and the infinite cluster]
Consider Bernoulli bond percolation on $\mathbb{Z}^d$: to each edge $e$ of the lattice,
independently assign $X_e \sim \mathrm{Bernoulli}(p)$ (edge \emph{open} if $X_e = 1$).
Let $C_\infty$ be the event that the origin belongs to an infinite open cluster.
\begin{enumerate}
  \item[(a)] Show that $C_\infty$ is a tail event with respect to the family $(X_e)_{e}$
             for any enumeration of the edges.
  \item[(b)] Conclude that $P_p(C_\infty) \in \{0, 1\}$.
  \item[(c)] For $d \ge 2$, there exists a critical probability $p_c \in (0,1)$ such that
             $P_p(C_\infty) = 0$ for $p < p_c$ and $P_p(C_\infty) = 1$ for $p > p_c$.
             Briefly explain why this dichotomy is consistent with (b).
\end{enumerate}
\end{exercise}

\begin{solution}
\textbf{(a)} The event $C_\infty$ is determined by the existence of an infinite self-avoiding
path starting from the origin.  Whether such a path exists is not changed by modifying
any finite collection of edges, because we can always reroute finitely many steps.
Formally, for any finite set of edges $F$, the occurrence of $C_\infty$ coincides (a.s.)
with the occurrence of the same event after removing $F$ from consideration.
Hence $C_\infty \in \sigma(X_e : e \notin F)$ for every finite $F$, which means
$C_\infty \in \mathcal{T} = \bigcap_F \sigma(X_e : e \notin F)$.

\textbf{(b)} Since $(X_e)$ are i.i.d., Kolmogorov's 0-1 Law applies and
$P_p(C_\infty) \in \{0,1\}$.

\textbf{(c)} The function $p \mapsto P_p(C_\infty)$ is non-decreasing (more open edges
can only help connectivity).  By (b) it can only take values in $\{0,1\}$, so it
must jump at some threshold $p_c$: below $p_c$ the probability is $0$, above it is $1$.
\end{solution}

% ----------------------------------------------------------------
\begin{exercise}[Symmetric nondegenerate random walk oscillates]
\label{ex:oscillate}
Let $X_1, X_2, \ldots$ be i.i.d.\ with distribution $\mu$ that is symmetric about $0$
(i.e., $\mu = \mu \circ (-\mathrm{id})^{-1}$) and nondegenerate (i.e., $P(X_1 = 0) < 1$).
Let $S_n = X_1 + \cdots + X_n$.  Show that
\[
  P\!\Bigl(\limsup_{n\to\infty} S_n = +\infty\Bigr) = 1
  \quad\text{and}\quad
  P\!\Bigl(\liminf_{n\to\infty} S_n = -\infty\Bigr) = 1.
\]
In particular, we are in case~(4) of the Four Possibilities Theorem for random walks on
$\mathbb{R}$ (i.e., $\limsup S_n = +\infty$ and $\liminf S_n = -\infty$ almost surely).
\end{exercise}

\begin{solution}
By the Hewitt--Savage 0-1 Law, $\limsup_n S_n$ is a.s.\ equal to a constant
$c \in [-\infty, +\infty]$ (it is a permutation-invariant measurable function of
$(X_1, X_2, \ldots)$, hence $0$ or $\pm\infty$).

Suppose for contradiction that $\limsup_n S_n = c < +\infty$ a.s. Since $\mu$ is
nondegenerate and symmetric, there exists $\epsilon > 0$ with $P(X_1 > \epsilon) = P(X_1 < -\epsilon) > 0$.
Then $P(X_1 > \epsilon) > 0$ implies $S_n \to +\infty$ with positive probability on certain
realizations (by the law of large numbers applied after conditioning), contradicting
finiteness of $c$.

More precisely: by symmetry, $(-S_n)$ has the same distribution as $(S_n)$, so
$\limsup(-S_n) = -\liminf S_n$ must equal the same constant $c$.
If $\limsup S_n = c$ a.s., then $\liminf S_n = -c$ a.s. For $c$ to be finite and
$-c$ to also be finite, we would need $S_n$ to stay in $[-c, c]$ forever, which is
impossible when $P(X_1 \ne 0) > 0$ (the walk escapes any bounded interval a.s.\ by
recurrence of returns to $0$ and the strong Markov property).  Hence $c = +\infty$.

By the symmetry argument applied to $-S_n$, we also have $\liminf S_n = -\infty$
a.s., placing us in case (4). \hfill$\square$
\end{solution}

\begin{remark}
This exercise provides an alternative proof that the simple symmetric random walk on
$\mathbb{Z}$ is recurrent: the walk oscillates between $+\infty$ and $-\infty$, so it
must return to $0$ infinitely often.
\end{remark}

\end{document}
